\clearpage
\VerdierPhysicalPageAnchor{0194}%
\begingroup\small\setlength{\parskip}{0pt}\setlength{\emergencystretch}{1.5em}
\setlength{\abovedisplayskip}{4pt plus 2pt minus 2pt}%
\setlength{\belowdisplayskip}{4pt plus 2pt minus 2pt}%

\VerdierSemanticAnchor{ch3d-2-2-6-proof-cone-cont}
(l'objet $L_0$ est le composant de degré $0$ de $c(m)$). La cohomologie de
$c(m)$ est nulle en degré différent de $-1$, et égale à $X$ en degré
$-1$. Le morphisme $L_n\longrightarrow L_n\oplus Z_{n-1}$ dont la
matrice est de la forme
$\begin{pmatrix}\operatorname{id}\\ *\end{pmatrix}$ admet une rétraction.
Il existe donc un morphisme
$v\colon Z_{n-1}\longrightarrow L_{n-1}\oplus L'_{n-2}$ tel que le
complexe :

\VerdierSemanticAnchor{ch3d-2-2-6-reduced-cone}
\[
\begin{aligned}
\cdots\longrightarrow0\longrightarrow0\longrightarrow Z_{n-1}
\xrightarrow{v}L_{n-1}\oplus L'_{n-2}\longrightarrow\cdots
\longrightarrow L_1\oplus L'_0\\
\longrightarrow L_0\longrightarrow0\longrightarrow\cdots
\end{aligned}
\]
soit quasi-isomorphe au complexe $c(m)$. D'autre part, le noyau du
morphisme $L_1\oplus L'_0\longrightarrow L_0$ est, d'après la propriété
$(E_4)$, un élément de $M$. Désignons par $L''_1$ ce noyau. Le complexe :

\VerdierSemanticAnchor{ch3d-2-2-6-resolution-x}
\[
\begin{aligned}
0\longrightarrow0\longrightarrow Z_{n-1}
\longrightarrow L_{n-1}\oplus L'_{n-2}\longrightarrow\cdots
\longrightarrow L_2\oplus L'_1\\
\longrightarrow L''_1\longrightarrow X\longrightarrow0\longrightarrow\cdots
\end{aligned}
\]
est donc acyclique. On a ainsi construit une résolution de longueur $n-1$
de l'objet $X$ et, par suite, $M\text{-}\dim(X)\leq n-1$. Par ailleurs,
il est clair que $M\text{-}\dim(X)\geq M\text{-}\dim(Y)-1$.

\VerdierSemanticAnchor{ch3d-2-2-6-conclusion}
Ceci achève la démonstration du lemme (2.2.6) et de la proposition (2.2.7).

\VerdierSemanticAnchor{ch3d-2-3}
\medskip\noindent\textbf{2.3. Quelques lemmes sur les résolutions.}

\VerdierSemanticAnchor{ch3d-2-3-1}
\medskip\noindent\textbf{Proposition 2.3.1.} \textbf{[22]}\quad
\textit{a) Soit $M$ un ensemble d'objets de $\mathcal A$ possédant les
propriétés $(E_1)$, $(E_2)$ et $(E_3)$ (2.2.1). Soient
$Y\in\operatorname{Ob}\operatorname{comp}^{-}(\mathcal A)$ un complexe
dont les composants appartiennent à $M^\wedge$ (2.2.1) et
$m\colon X\longrightarrow Y$ un quasi-isomorphisme. Il existe un complexe
$L^\bullet\in\operatorname{Ob}\operatorname{comp}^{-}(\mathcal A)$ dont
les composants appartiennent à $M$ et un quasi-isomorphisme
$w\colon L^\bullet\longrightarrow X$.}

\noindent\textit{b) Réciproquement, soit $M$ un ensemble d'objets de
$\mathcal A$ possédant la propriété $(E_1)$. Supposons que pour tout
complexe $Y\in\operatorname{Ob}\operatorname{comp}^{-}(\mathcal A)$ dont
les composants sont des éléments de $M$, et pour tout quasi-isomorphisme
$m\colon X\longrightarrow Y$, il existe un complexe $L^\bullet$ dont les
composants appartiennent à $M$ et un quasi-isomorphisme
$w\colon L^\bullet\longrightarrow X$. Alors l'ensemble $M$ possède les
propriétés $(E_2)$ et $(E_3)$.}

\VerdierSemanticAnchor{ch3d-2-3-1-proof}
Démontrons a). Tout d'abord, il existe un complexe
$X'\in\operatorname{Ob}\operatorname{comp}^{-}(\mathcal A)$ et un
quasi-isomorphisme $X'\longrightarrow X$ (1.2.1). On peut donc se ramener
au cas où $X$ est un objet de $\operatorname{comp}^{-}(\mathcal A)$.
D'après (2.1.4) et (2.1.5), il existe une suite de morphismes :

\VerdierSemanticAnchor{ch3d-2-3-1-xn-sequence}
\[
u_n\colon X_n\longrightarrow X_{n-1},
\qquad n\in\mathbf Z,
\]
\endgroup
